\documentclass[10pt,a4paper]{article}
\usepackage[margin=18mm]{geometry}
\usepackage{enumitem}
\usepackage[hidelinks]{hyperref}
\usepackage{kotex}
\setlength{\parindent}{0pt}
\begin{document}
{\LARGE 조형곤}\\[3pt]
건물에너지 연구자 · 과학 소프트웨어 개발자\\
서울대학교 건물시뮬레이션연구실\\
\href{mailto:hyeonggon.jo@snu.ac.kr}{hyeonggon.jo@snu.ac.kr} \quad \url{https://github.com/Gonie-Gonie} \quad \url{https://orcid.org/0009-0004-8821-275X}

\section*{소개}
건물에너지 성능, 그린리모델링 평가, HVAC 제어, 시뮬레이션 방법론을 연구합니다. 공학 연구를 반복·공유하기 쉽게 만드는 작고 검사 가능한 도구도 직접 개발합니다.

\section*{학력}
\begin{itemize}[leftmargin=*,nosep]
\item \textbf{건축학과 공학박사}, 서울대학교 \hfill 2022.03–2026.02\\\emph{설계단계 의사결정 지원을 위한 동적 시뮬레이션 도구 단순화 방법}
\item \textbf{건축학과 공학석사}, 서울대학교 \hfill 2020.03–2022.02\\\emph{확률론적 주광 모델 및 전기조명 제어}
\item \textbf{건축학과 공학사}, 서울대학교 \hfill 2016.03–2020.02
\end{itemize}

\section*{주요 연구 프로젝트}
\begin{itemize}[leftmargin=*]
\item \textbf{그린리모델링 건축물 성능평가} \hfill 2025.06–2025.10\\한국부동산원; 연구원 · 시뮬레이션 및 데이터 분석. 완료된 그린리모델링 사업의 운영 데이터를 수집하고 동적 시뮬레이션을 활용해 실현 절감량과 성능 격차를 평가했습니다.
\item \textbf{기존 건축물 에너지·온실가스 시뮬레이터 프로토타입} \hfill 2024.07–2025.12\\국토안전관리원; 시뮬레이션 구조 및 EnergyPlus 워크플로. 구조화된 스프레드시트 입력을 반복 가능한 EnergyPlus 모델과 리모델링 대안으로 변환하는 프로토타입을 개발했습니다.
\item \textbf{그린리모델링 온실가스 감축계수 데이터베이스} \hfill 2022.12–2024.12\\국토안전관리원; EnergyPlus 모델링, 대규모 시뮬레이션, 의사결정 도구 개발. 노후 건축물 아키타입을 대상으로 대규모 시뮬레이션을 수행해 감축계수와 장기 배출량 추정체계를 구축했습니다.
\item \textbf{산업 FAB 유틸리티 시스템 통합 최적화 시뮬레이터} \hfill 2022.05–2025.03\\삼성물산; 시스템 모델링, 상태추정, 제어 최적화, Python 시뮬레이터. 공조기, 열회수, 열원 시스템의 물리·데이터 기반 모델을 개발하고 노드 기반 시뮬레이터로 연결했습니다.
\end{itemize}

\section*{주요 논문}
\begin{enumerate}[leftmargin=*]
\item Park, C. H., Lee, S. J., Jo, H. G., Park, C. S. (2026). Implications of geometry simplification in dynamic building energy modeling: Comparative analysis of detailed and simplified geometries. Energy and Buildings, 117839.
\item Kim, J. H., Kim, Y. S., Jo, H. G., Urabe, E., Mun, J., Shin, Y., Park, Y., Park, C. S. (2025). Accuracy, generalizability, and extrapolation ability of physics-based, data-driven, and hybrid models for real-life cooling towers. Building and Environment, 112756.
\item Mun, J., Jo, H. G., Kim, J. H., Urabe, E., Mun, J., Park, Y., Park, C. S. (2025). Causality-driven, interpretable surrogate modeling for DOAS preheating control. Energy and Buildings, 116625.
\item Kim, J. H., Kim, Y. S., Jo, H. G., Mun, J., Park, C. S. (2025). Domain-invariant representation learning for generalizable chiller model: a real-world case study. Energy and Buildings, 116168.
\item Ra, S. J., Jo, H. G., Park, C. S. (2025). Real-time model predictive control of energy recovery ventilators for a school building. Science and Technology for the Built Environment, Vol. 31, No. 9, pp. 1155-1166. https://doi.org/10.1080/23744731.2025.2539048
\item Jo, H. G., Park, C. H., Lee, S. J., Mun, H. J., Park, C. S. (2025). Running EnergyPlus Using Spreadsheet Inputs: Focus on Geometry Information. Proceedings of the Fall Conference of the Architectural Institute of Korea, Vol. 45, No. 2, pp. 574-575. October 29-31, 2025, Science and Technology Convention Center, Seoul.
\item Mun, J., Jo, H. G., Park, C. S. (2025). Toward scalable prediction of indoor thermal dynamics: Neural-network-implemented state-space (NNiSS) model. Energy and Buildings 331, 115359.
\item Jo, H. G., Kwon, S. O., Kim, U. H., Yoo, Y. S., Park, C. S. (2024). Establishment and Application of a Greenhouse Gas Reduction Factor Database for Policy Decision-Making in Green Remodeling Projects. Proceedings of the Fall Conference of the Architectural Institute of Korea, Vol. 44, No. 2, pp. 501-502. October 23-25, 2024, The-K Hotel Gyeongju, Gyeongju, Gyeongsangbuk-do.
\item Jo, H. G., Kim, Y. S., Kim, J. H., Park, C. S., Urabe, E., Mun, J., Shin, Y., Park, Y. (2024). Optimal Control Strategies of a Heat Recovery System for a Building. Proceedings of SASBE 2024, November 7-9, 2024, Auckland, New Zealand. https://doi.org/10.1007/978-981-96-4051-5
\item Jo, H. G., Lee, S. J., Park, C. H., Mun, H. J., Park, C. S. (2024). Running EnergyPlus Using Spreadsheet Inputs: Focus on HVAC Systems. Proceedings of the Fall Conference of the Architectural Institute of Korea, Vol. 45, No. 2, pp. 576-577. October 29-31, 2025, Science and Technology Convention Center, Seoul.
\item Lee, S. J., Jo, H. G., Yoo, Y. S., Park, C. H., Park, C. S. (2023). Issues in Quasi-Steady-State Building Energy Analysis Tools: Focusing on ECO2, the Korean Building Energy Efficiency Rating Certification Evaluation Program. Journal of the Architectural Institute of Korea, Vol. 40, No. 1, pp. 169-176.
\item Jo, H. G., Choi, S. H., Park, C. S. (2023). Linear regression based indoor daylight illuminance estimation with simple measurements based on daylight coefficient method. Journal of Building Performance Simulation 16(5), pp. 574-587. https://doi.org/10.1080/19401493.2023.2185684
\item Jo, H. G., Choi, S. H., Heo, S. Y., Park, C. S. (2022). Calibration Method for a Daylight Simulation Model Using Field Measurement Data. Proceedings of the Fall Conference of the Korean Institute of Architectural Sustainable Environment and Building Systems, pp. 119-121. November 25, 2022, Yonsei University, Seoul.
\item Jo, H. G., Choi, S. H., Park, C. S. (2022). Indoor Daylight Estimation Model Based on the Daylight Coefficient Concept. Proceedings of the Spring Conference of the Architectural Institute of Korea, Vol. 42, No. 1, pp. 399-400. April 28-29, 2022, The-K Hotel, Seoul.
\item Jo, H. G., Choi, S. H., Park, C. S. (2022). Real-time indoor daylight illuminance simulation of an existing building using minimal data. Proceedings of the 2022 Winter Simulation Conference, December 11-14, 2022, Singapore, Singapore.
\item Jo, H. G., Kim, Y. S., Park, C. S. (2021). Application of daylighting Gaussian process emulator to lighting control of an existing building. Proceedings of the 2021 Winter Simulation Conference, December 13-17, 2021, Phoenix, Arizona, USA.
\item Jo, H. G., Ra, S. J., Kim, J. H., Park, C. S. (2021). Real-Time Electric Lighting Control Considering the Uncertainty of a Daylight Model. Proceedings of the Fall Conference of the Architectural Institute of Korea, Vol. 41, No. 2, pp. 332-333. October 27-30, 2021, Yeosu Expo Convention Center, Yeosu, Jeollanam-do.
\item Jo, H. G., Kim, Y. S., Park, C. S. (2021). Stochastic Daylight Model and Its Application to Model Predictive Control in Large-Space Buildings. Journal of the Architectural Institute of Korea, Vol. 37, No. 12, pp. 255-264.
\item Jo, H. G., Park, C. S. (2020). Illuminance prediction using a hybrid simulation model based on a single reference measurement. Proceedings of the 2020 Winter Simulation Conference, December 14-18, 2020, Orlando, USA. Virtual Conference.
\item Jo, H. G., Shin, H. S., Kim, Y. S., Park, C. S. (2020). Prediction of Illuminance Ratios Between Indoor Points Using Gaussian Processes. Proceedings of the Fall Conference of the Korean Institute of Architectural Sustainable Environment and Building Systems, pp. 153-154. November 19-21, 2020, Korea National University of Transportation, Chungju.
\end{enumerate}

\section*{연구 소프트웨어}
\begin{itemize}[leftmargin=*]
\item \textbf{EPlusSimple} --- 스프레드시트 입력을 검증하고 EnergyPlus 모델 생성·실행·결과 요약까지 연결하는 워크플로입니다.\\\url{https://github.com/snu-bslab/EPlusSimple}
\item \textbf{semantic-idf} --- EnergyPlus IDF·epJSON 모델을 검사·진단·정리·비교·변환하는 데스크톱 워크벤치입니다.\\\url{https://github.com/Gonie-Gonie/semantic-idf}
\item \textbf{oo-docs} --- 하나의 구조화된 Python 문서 트리를 DOCX·PDF·HTML로 렌더링하며 데이터·그림·인용·레이아웃을 공유합니다.\\\url{https://github.com/Gonie-Gonie/oo-docs}
\item \textbf{eng-lang} --- 단위·스키마·축·출처·검토 산출물을 보존하는 실험적 공학 워크플로 언어입니다.\\\url{https://github.com/Gonie-Gonie/eng-lang}
\item \textbf{rusted-energyplus} --- EnergyPlus 알고리즘과 호환 동작을 근거 중심으로 재구현하는 Rust 실험 프로젝트입니다.\\\url{https://github.com/Gonie-Gonie/rusted-energyplus}
\item \textbf{hvac-studio} --- 상호연결 건물설비 모델을 구성하는 컴포넌트·노드 기반 저작환경과 런타임입니다.\\\url{https://github.com/Gonie-Gonie/hvac-studio}
\end{itemize}

\section*{주요 수상}
\begin{itemize}[leftmargin=*]
\item 2025-12-02 --- \textbf{ISHVAC 2025 Best Paper Award}, 제14회 International Symposium on Heating, Ventilating and Air-Conditioning. Causality-driven, interpretable surrogate modelling for DOAS preheating control
\item 2024-10-25 --- \textbf{우수발표논문상}, 대한건축학회. 그린리모델링 정책 의사결정을 위한 온실가스 감축계수 데이터베이스 구축 및 적용
\end{itemize}
\end{document}
